\section{Программная реализация}

В соответствии с архитектурой, описанной в разделе~3, программная часть проекта была реализована в виде двух независимо собираемых частей: серверного приложения на платформе ASP.NET~Core и клиентского одностраничного приложения на библиотеке React.
Сборка серверного приложения выполняется командой \code{dotnet build}, клиентского~--~командой \code{npm run build}.
База данных управляется через инструмент Liquibase, миграции применяются автоматически при первом запуске контейнера.

\subsection{Структура серверной части}

Серверная часть приложения сгруппирована по подкаталогам в соответствии с назначением входящих в них классов.

Каталог \code{Controllers} содержит REST-контроллеры, по одному на каждую группу сущностей.
Каждый контроллер наследуется от \code{ControllerBase}, помечен атрибутом \code{[ApiController]} и атрибутом маршрутизации \code{[Route("api/v1/...")]}.
Для простых CRUD-операций контроллер делегирует выполнение соответствующему репозиторию напрямую.
Для составных операций~--~таких как запрос на удаление абитуриента, подтверждение или отклонение запроса, валидация данных~--~контроллер выступает оркестратором: координирует обращения к нескольким репозиториям, извлекает идентификатор текущего пользователя из утверждений JWT и при необходимости инициирует пересчёт результатов отбора через \code{SelectionService}.

Каталог \code{Repositories} содержит репозитории, по одному на каждую таблицу базы данных.
Репозитории работают с базой данных через драйвер Npgsql, выполняя параметризованные SQL-запросы.
Каждый репозиторий поддерживает базовый набор операций CRUD, а также специфичные для предметной области операции (например, \code{GetByFacultyAsync}, \code{GetPagedAsync}).

Каталог \code{Services} содержит сервисы, реализующие бизнес-логику, выходящую за рамки CRUD: \code{JwtService} для генерации и валидации JSON Web Tokens; \code{SelectionService} для пересчёта результатов отбора; \code{ExcelExportService} для формирования файлов Excel; \code{AuditLogger} для записи журнала действий пользователей; \code{AuthLogger} для записи журнала попыток входа; \code{CurrentUserService} для получения сведений о текущем пользователе из утверждений (claims) JWT.

Каталог \code{Dtos} содержит объекты передачи данных, разделённые на подкаталоги \code{Requests} (входящие данные) и \code{Responses} (исходящие данные).
Для постраничной выдачи используется обобщённый тип \code{PagedResponse<T>}.

Каталог \code{Constants} содержит вспомогательные классы констант, в том числе \code{UserRoles} с предопределёнными строками для атрибута авторизации.

Каталог \code{Resources} содержит маркерный класс \code{SharedResources} и файлы ресурсов \code{SharedResources.ru.resx} и \code{SharedResources.en.resx} с переводами сообщений об ошибках.

\subsection{Регистрация сервисов и конфигурация приложения}

Все сервисы и репозитории регистрируются в контейнере зависимостей в файле \code{Program.cs}.
Файл также настраивает: подключение к базе данных PostgreSQL через строку соединения из конфигурации; аутентификацию по JWT с алгоритмом HS256 и валидацией издателя, времени жизни и подписи токена; авторизацию с предопределёнными политиками для каждой роли; локализацию с поддержкой языков \code{ru} (по умолчанию) и \code{en}; CORS-политику, разрешающую запросы от настроенных в конфигурации источников с обязательным экспонированием заголовка \code{Content-Disposition} (необходим для скачивания файлов Excel).

При первом запуске приложения, если в базе данных не существует ни одного пользователя, выполняется создание учётной записи администратора с именем пользователя и паролем, заданными в переменных окружения \code{InitialAdmin\_\_Username} и \code{InitialAdmin\_\_Password} (значения по умолчанию~--~\code{admin} и \code{Admin@123}).
Учётной записи присваивается роль \code{SuperAdmin}.

\subsection{Реализация аутентификации и авторизации}

Аутентификация реализована в контроллере \code{AuthController}, эндпоинт \code{POST /api/v1/auth/login}.
Контроллер принимает имя пользователя и пароль, выполняет поиск учётной записи в базе данных, проверяет хеш пароля и, в случае успеха, генерирует JSON Web Token через сервис \code{JwtService}.
Токен содержит утверждения с идентификатором пользователя, именем пользователя и ролью; срок действия токена~--~8 часов.
Параллельно с генерацией токена сервис \code{AuthLogger} записывает успешную попытку входа в таблицу \code{auth\_log}.
Полный текст сервиса \code{JwtService.cs} приведён в приложении~А.
При неуспешной попытке (несуществующий пользователь, неверный пароль, заблокированная учётная запись) также делается запись в журнал с соответствующим статусом.

Авторизация реализована декларативно через атрибуты на уровне методов контроллеров.
Для удобства использования в классе \code{UserRoles} определены строковые константы, соответствующие группам ролей с одинаковыми правами: \code{SuperAdmin}, \code{WriteStructure} (SuperAdmin плюс DataAdministrator), \code{WriteApplicants} (SuperAdmin, DataAdministrator, Auditor, AdmissionsOperator), \code{WriteAudit} (SuperAdmin, DataAdministrator, Auditor).
Например, метод создания факультета помечен атрибутом \code{[Authorize(Roles = UserRoles.WriteStructure)]}.

\subsection{Реализация алгоритма пересчёта результатов отбора}

Алгоритм пересчёта результатов отбора реализован в сервисе \code{SelectionService} с разделением на координирующий метод \code{RecalculateAllAsync} и чистую in-memory логику в статическом классе \code{FullRecalculationCore}.
Такое разделение позволяет проверить алгоритм юнит-тестами на готовом снимке без подключения к базе данных.

Метод \code{RecalculateAllAsync} выполняет последовательность действий: открывает транзакцию к базе данных через \code{NpgsqlConnection.BeginTransactionAsync}; захватывает глобальную транзакционную блокировку через функцию \code{pg\_advisory\_xact\_lock(0)}, что исключает параллельный пересчёт при одновременном обновлении данных несколькими пользователями; загружает снимок базы данных шестью пакетными SQL-запросами (по одному на тип сущности~--~заявки, категории, планы, критерии групп, оценки, типы критериев), исключая итеративное обращение к базе данных в цикле по записям; передаёт снимок в метод \code{FullRecalculationCore.Run}; полностью очищает таблицу \code{selectedapplicants}; массово вставляет результат через операцию \code{COPY BINARY} для достижения максимальной производительности; записывает в журнал аудита событие \code{recalculate\_all}; подтверждает транзакцию.

Для массовой вставки результата использован интерфейс \code{NpgsqlBinaryImporter}, экспортируемый методом \code{BeginBinaryImportAsync} драйвера Npgsql; формат вставки~--~бинарный, что устраняет накладные расходы на разбор текстовых представлений значений на стороне сервера базы данных.
Полный текст модуля \code{SelectionService.cs} приведён в приложении~А.

Метод \code{FullRecalculationCore.Run} реализует алгоритм очереди с вытеснением, описанный в разделе~3.7.
Метод чист: на одинаковом входе всегда производит одинаковый выход, что упрощает написание юнит-тестов.
Полный текст модуля \code{FullRecalculationCore.cs} приведён в приложении~А.
Сравнение кандидатов в пуле категории выполняется компаратором \code{CategoryComparer}, реализующим интерфейс \code{IComparer<Candidate>}.
Компаратор спроектирован так, чтобы максимум \code{SortedSet} соответствовал наиболее слабому кандидату~--~это упрощает операцию вытеснения, выполняемую при превышении квоты.
Полный текст модуля \code{CategoryComparer.cs} приведён в приложении~А.

Для критериев с направлением оптимизации \code{lower\_is\_better} значение умножается на минус единицу при построении вектора оценок, что позволяет использовать единый компаратор по принципу <<больше~--~лучше>> без ветвления внутри сравнения.

\subsection{Реализация подсистемы аудита}

Сервис \code{AuditLogger} реализует интерфейс \code{IAuditLogger}, состав методов которого приведён в таблице~\ref{tab:auditlogger}.

\begin{longtblr}[
  caption = {Методы интерфейса IAuditLogger},
  label = {tab:auditlogger},
]{
  colspec = {|Q[l,wd=6.5cm]|X|},
  width = \textwidth,
  hlines, vlines,
  rowhead = 1,
  row{1} = {halign=c},
}
Метод & Назначение \\
\code{LogCreateAsync} & Запись действия \code{create} с состоянием после изменения \\
\code{LogUpdateAsync} & Запись действия \code{update} с разностью состояний до и после \\
\code{LogDeleteAsync} & Запись действия \code{delete} с состоянием до изменения \\
\code{LogValidateAsync} & Запись действия \code{validate} (ручное подтверждение валидации) \\
\code{LogInvalidateAsync} & Запись действия \code{invalidate} с пометкой ручного или автоматического сброса \\
\code{LogDeleteConfirmedAsync} & Запись действия \code{delete\_confirmed} \\
\code{LogDeleteRejectedAsync} & Запись действия \code{delete\_rejected} \\
\code{LogRecalculationAsync} & Запись действия \code{recalculate\_all} с числом зачисленных \\
\code{ResetValidationIfNeededAsync} & Условный сброс поля \code{validated} при изменении данных абитуриента \\
\end{longtblr}

Сервис внедряется во все репозитории и сервисы, выполняющие изменения данных, через конструктор.
Внутренний метод \code{WriteAsync} формирует параметризованный SQL-запрос \code{INSERT INTO audit\_log} с приведением поля \code{changes} к типу \code{jsonb} непосредственно в строке запроса, что устраняет необходимость преобразования параметра на стороне приложения.
Идентификатор и имя пользователя извлекаются из сервиса \code{ICurrentUserService}, который читает соответствующие \code{Claims} из текущего \code{HttpContext}.

Метод \code{ResetValidationIfNeededAsync} проверяет текущее значение поля \code{validated} абитуриента и выполняет обновление только в случае, когда оно равно \code{true}, что исключает избыточные записи в журнал при изменении данных уже невалидированного абитуриента.
При сбросе автоматически записывается событие \code{invalidate} с пометкой \code{automatic = true}, отделяющей автоматический сброс от явного снятия валидации пользователем.

Для просмотра журнала действий предусмотрен контроллер \code{AuditController} с эндпоинтом \code{GET /audit/log}, поддерживающим фильтрацию по дате, действию, типу сущности и пользователю, а также эндпоинт \code{GET /audit/entity-history}, возвращающий полную историю изменений конкретной сущности.
Полный текст модуля \code{AuditLogger.cs} приведён в приложении~А.

\subsection{Реализация логического удаления абитуриентов}

Двухступенчатое удаление реализовано через таблицу \code{applicant\_deletion\_requests}.
При вызове метода \code{DELETE /applicants/\{id\}} репозиторий создаёт в таблице запись со статусом \code{pending}, связывая её с идентификатором текущего пользователя и идентификатором абитуриента.
После этого вызывается \code{SelectionService.RecalculateAllAsync}, который при формировании снимка исключает абитуриентов с активными запросами на удаление.

Для подтверждения и отклонения запросов реализованы эндпоинты \code{POST /audit/pending-deletions/\{id\}/confirm} и \code{POST /audit/pending-deletions/\{id\}/reject}, доступные пользователям с правом \code{WriteAudit}.
При подтверждении статус запроса меняется на \code{confirmed}, и абитуриент остаётся скрытым из всех выборок.
При отклонении статус меняется на \code{rejected}, флаг валидации абитуриента сбрасывается в значение \code{false}, и запускается повторный пересчёт результатов отбора с включением вернувшегося абитуриента.
После отклонения возможно создание нового запроса на удаление того же абитуриента.

\subsection{Структура клиентской части}

Клиентская часть структурирована по принципу разделения по типу содержимого.

Каталог \code{src/api} содержит модули, инкапсулирующие вызовы REST API.
Каждый модуль соответствует одной группе эндпоинтов сервера: \code{authApi.js}, \code{applicantsApi.js}, \code{facultyApi.js} и так далее.
Все модули используют единый экземпляр Axios, настроенный в файле \code{src/api/index.js}.

Каталог \code{src/pages} содержит компоненты страниц, отображаемых по разным маршрутам: \code{Login.jsx}, \code{Applicants.jsx}, \code{Faculties.jsx} и другие.
Маршрутизация описана в файле \code{src/AppRouter.jsx} с использованием библиотеки React~Router.

Каталог \code{src/components} содержит переиспользуемые компоненты: \code{ContentHeader} с переключателем языка, \code{ProtectedRoute} для защиты маршрутов по роли, общие компоненты таблиц и форм.

Каталог \code{src/contexts} содержит провайдеры глобальных контекстов: \code{AuthContext}, \code{ThemeContext}, \code{LanguageContext}.

Каталог \code{src/hooks} содержит пользовательские хуки, в том числе \code{usePermissions}, возвращающий булевые флаги для проверки прав пользователя на запись.

Каталог \code{src/i18n} содержит конфигурацию библиотеки react-i18next и файлы переводов \code{locales/ru/translation.json}, \code{locales/en/translation.json}.

\subsection{Интерцептор Axios}

Интерцептор запроса, настроенный на едином экземпляре Axios, прикрепляет к каждому исходящему запросу заголовок \code{Authorization: Bearer <token>}, считывая токен из локального хранилища браузера по ключу \code{auth}.
Также интерцептор добавляет заголовок \code{Accept-Language} со значением языка из локального хранилища по ключу \code{language}~--~это позволяет серверу выбирать нужный язык сообщений об ошибках.
Чтение языка из локального хранилища напрямую (без использования React-хуков) обусловлено тем, что хуки React недоступны в интерцепторах Axios.
Полный текст модуля \code{api/index.js} с реализацией интерцепторов приведён в приложении~А.

Интерцептор ответа обрабатывает HTTP-статус 401 (Unauthorized): при его получении интерцептор очищает локальное хранилище и перенаправляет пользователя на страницу входа.
Это обеспечивает корректное поведение приложения при истечении срока действия токена.

\subsection{Хук серверной пагинации таблиц}

Все таблицы клиентской части обязаны использовать серверную пагинацию и серверные фильтры~--~это требование к проекту обусловлено невозможностью загружать в браузер десятки тысяч записей абитуриентов целиком.
Для централизации этого поведения разработан пользовательский хук \code{useServerTable}, инкапсулирующий состояние пагинации, фильтров, сортировки и загрузки.

Хук принимает на вход стабильную (обёрнутую в \code{useCallback}) функцию запроса данных и возвращает объект с текущими значениями состояния, обработчиком события изменения таблицы и подготовленным объектом параметра \code{pagination} для компонента \code{Table} библиотеки Ant Design.
При изменении страницы, размера страницы, фильтров или сортировки хук автоматически запускает повторный запрос с актуальными параметрами.
Применение хука обеспечивает единообразный пользовательский опыт во всех табличных страницах приложения и устраняет потребность в дублировании кода управления состоянием таблиц.
Полный текст модуля \code{useServerTable.js} приведён в приложении~А.

\subsection{Локализация сообщений об ошибках}

Сообщения об ошибках, формируемые серверной частью, локализуются через встроенный в платформу~.NET механизм \code{Microsoft.Extensions.Localization}.
В каталоге \code{Resources} расположен маркерный класс \code{SharedResources} и два файла ресурсов: \code{SharedResources.ru.resx} с русскими переводами (по умолчанию) и \code{SharedResources.en.resx} с английскими.
Ключи ресурсов записываются по соглашению \code{Domain.Context} с использованием PascalCase, например \code{Faculty.Name.Required}, \code{Applicant.NotFound} или \code{ApplicantEvaluationValue.ValueOutOfRange}.

В контроллерах локализатор внедряется через интерфейс \code{IStringLocalizer<SharedResources>}.
При построении объекта-ответа сообщение явно приводится к типу \code{string} для корректной сериализации в JSON: \code{return NotFound(new \{ message = (string)\_localizer["Applicant.NotFound", id] \});}.
Все сообщения об ошибках возвращаются клиенту в едином формате \code{\{ "message": "..." \}}, что позволяет на клиенте читать значение через выражение \code{err.response?.data?.message} независимо от типа ошибки.

В DTO-классах входящих запросов атрибуты валидации \code{[Required]}, \code{[StringLength]} и аналогичные принимают ключи ресурсов в свойстве \code{ErrorMessage}, а не локализованные строки.
Маршрутизация этих ключей через \code{SharedResources} обеспечивается настройкой \code{AddDataAnnotationsLocalization} в \code{Program.cs}, что позволяет использовать общую систему ключей и для бизнес-логики контроллеров, и для атрибутной валидации.

Выбор языка ответа выполняется на основе заголовка \code{Accept-Language}, передаваемого клиентом в каждом запросе.
Поддерживаются культуры \code{ru} (по умолчанию) и \code{en}; промежуточное программное обеспечение \code{RequestLocalization} устанавливает текущую культуру потока обработки запроса в соответствии с этим заголовком.

\subsection{Управление миграциями базы данных}

Все изменения схемы базы данных описаны в виде набора changeset-ов в файле \code{db.changelog-master.yaml} с подключёнными SQL-сценариями из каталога \code{liquibase/sql}.
На момент окончания дипломного проектирования в проекте применены 27 миграций, реализующих эволюцию схемы от первоначальной версии до текущей.
Запуск Liquibase осуществляется через отдельный контейнер Docker, описанный в файле \code{liquibase/docker-compose.yml}.

\subsection{Развёртывание системы}

Развёртывание системы построено на основе контейнеризации через платформу Docker \citesrc{docker}.
Контейнеризация обеспечивает воспроизводимость окружения от рабочей станции разработчика до промышленной эксплуатации, упрощает обновление компонентов и устраняет зависимость от конкретной версии операционной системы и установленного на сервере прикладного программного обеспечения.

Диаграмма развёртывания системы представлена на рисунке~\ref{fig:deployment}.

\begin{figure}[htbp]
\centering
\includegraphics[width=\textwidth]{deployment}
\caption{Диаграмма развёртывания системы}
\label{fig:deployment}
\end{figure}

В составе развёртывания выделены два узла: рабочая станция пользователя и сервер приложений.

На рабочей станции пользователя выполняется только веб-браузер (поддерживаются последние версии Chrome, Firefox и Safari).
Браузер исполняет загруженную сборку клиентского одностраничного приложения и хранит JWT-токен доступа в локальном хранилище \code{localStorage}.
Установка какого-либо прикладного программного обеспечения на рабочую станцию не требуется, что упрощает развёртывание системы среди сотрудников приёмной комиссии.

Сервер приложений объединяет четыре программных компонента: веб-сервер статических файлов \code{nginx} (отдаёт сборку клиентского одностраничного приложения), серверную часть на платформе ASP.NET~Core с встроенным веб-сервером Kestrel (порт 5059), контейнер PostgreSQL 18 (порт 5432) и одноразовый контейнер Liquibase, применяющий миграции схемы при обновлении системы.
PostgreSQL и Liquibase развёртываются в Docker через отдельные файлы \code{infra/docker-compose.yml} и \code{liquibase/docker-compose.yml}.
Данные базы хранятся в именованном томе Docker \code{postgres\_data}, что обеспечивает их сохранность при пересоздании контейнера PostgreSQL.

Взаимодействие компонентов на сервере организовано следующим образом.
Веб-браузер пользователя обращается к веб-серверу статических файлов по протоколу HTTPS для загрузки сборки клиентского приложения, а затем к серверной части по REST API в формате JSON по протоколу HTTPS на порту~5059.
Серверная часть подключается к контейнеру PostgreSQL по протоколу TCP на порту~5432 через драйвер Npgsql.
Контейнер Liquibase подключается к тому же экземпляру PostgreSQL и выполняется однократно при запуске или обновлении системы~--~после применения недостающих миграций контейнер завершает работу.

Параметры подключения к базе данных (имя пользователя, пароль, имя базы, порт) задаются через файл переменных окружения \code{.env}, отдельный для каждого окружения (разработка, тестирование, эксплуатация).
Учётная запись администратора системы создаётся автоматически при первом запуске серверной части, если в таблице \code{users} нет ни одного пользователя; имя пользователя и пароль задаются переменными окружения \code{InitialAdmin\_\_Username} и \code{InitialAdmin\_\_Password} (значения по умолчанию~--~\code{admin} и \code{Admin@123}).

Сценарий обновления версии системы состоит из трёх шагов: остановки контейнера серверной части, запуска контейнера Liquibase для применения новых миграций схемы (контейнер завершается автоматически после успешного применения), запуска контейнера серверной части новой версии.
Контейнер PostgreSQL и его том данных при обновлении не пересоздаются, что обеспечивает сохранность данных приёмной кампании между версиями системы.
В случае отката версии порядок обратный, но в качестве дополнительной меры предосторожности необходимо создание резервной копии тома \code{postgres\_data} до применения миграций, поскольку не все миграции схемы обратимы.

Полные тексты конфигурационных файлов \code{docker-compose.yml} для PostgreSQL и Liquibase, \code{Dockerfile} контейнера Liquibase, шаблонов переменных окружения \code{.env.example} и пример SQL-миграции схемы приведены в приложении~А.

\subsection{Интернационализация интерфейса}

Поддержка двух языков интерфейса реализована на двух уровнях.

На клиентской части используется библиотека react-i18next с автоматическим определением языка через библиотеку i18next-browser-languagedetector.
Переключение языка осуществляется через компонент \code{Segmented} библиотеки Ant Design в заголовке страницы.
При смене языка интерфейс автоматически перерисовывается без перезагрузки страницы.
Локаль библиотеки Ant Design также переключается динамически между \code{ruRU} и \code{enUS} в файле \code{src/providers/GlobalProvider.jsx}.

На серверной части используется встроенный в платформу~.NET механизм \code{Microsoft.Extensions.Localization}.
Файлы ресурсов \code{SharedResources.ru.resx} и \code{SharedResources.en.resx} содержат переводы сообщений об ошибках.
Сообщения подключаются в контроллерах через инжектируемый интерфейс \code{IStringLocalizer<SharedResources>}.
Атрибуты валидации DTO используют ключи ресурсов в свойстве \code{ErrorMessage}, которые автоматически разрешаются через ту же систему ресурсов благодаря настройке \code{DataAnnotationsLocalization} в \code{Program.cs}.

Выбор языка ответа сервера выполняется на основе заголовка \code{Accept-Language}, передаваемого клиентом в каждом запросе.
По умолчанию используется русский язык.
